% P10-Evaluationsfeinschliff, 16.09.2026: Aussagegrenzen und Anforderungsbilanz; keine neuen Messwerte.
% P10-Core-Inhaltsprüfung, 15.09.2026: datierter Ergebnisanschluss; B-60 korrigiert.
% Belege: thesis-evidence/w2/core-content-audit-20260915/evaluation/
% P10-Methodenabgleich, 15.09.2026: Verfahren nach beanspruchter Eigenschaft begründet.
% Bestehende Prüffragen, Evidenzklassen, Ergebnisse und Reichweiten bleiben erhalten.
% P10-4b, 14.09.2026: Begriffs-/Beleganschlüsse und abgegrenzte Prüfaussagen präzisiert.
% M19 (12.09.2026): Leserführung und präzise Ergebnisdarstellung;
% Messwerte und historische Nachweisgrenzen bleiben erhalten.
% ============================================================
% §7.1 Evaluationsziele und Untersuchungsumfang — v01 (07.09.)
% NEUE 8er-Struktur (Schreibplan-Zeile 7.1: K4-Anforderungen als
% konkrete PRÜFFRAGEN [Evidenzbindung, kontrollierte Transformation/
% Mutation, Fortschreibung, Projektion, Durabilität]; was bleibt
% außerhalb). Budget ~0,8 S. + 1 Tabelle (Prüflandkarte — wird in
% 7.7 mit Befunden GESCHLOSSEN; ersetzt die alte Claim-to-Evidence-
% Matrix chap7.1table.tex inhaltlich in anforderungsbezogener Form).
% ============================================================
% HERKUNFT: Prüffragen = anforderungsbezogene Neuformung der
% Claim-to-Evidence-Matrix (alt chap7.1table.tex / E2E-EVIDENZ.md
% §0-Evidenzklassen) gemäß Schreibplan v2.4-Leitfrage; Eigenschaften
% = K4-Anforderungskategorien; Abschnittszuordnung = tex-v2-Struktur.
% ------------------------------------------------------------

\chapter{Evaluation}
\label{chap:evaluation}
% [KAPITEL-KOPF-FIX (08.09., nach PDF-30-Review R1): \chapter+\label
%  lagen im ALTEN 7.1 (kapitel-7/archive/chap7.1) und gingen beim
%  tex-v2-Umstieg verloren — PDF 30 zeigt deshalb die Evaluation als
%  6.10–6.17 OHNE Kapitel-7-Überschrift und "Kapitel ??" im Anhang.
%  Jetzt K6-Muster (wie chap6.1: \chapter + \label IN der ersten
%  Abschnittsdatei). ⚠ Overleaf-Hauptdatei prüfen: Zwischen dem
%  chap6.9- und diesem chap7.1-Include darf KEIN weiteres
%  \chapter{...} stehen (sonst doppeltes Kapitel). Nach dem Sync
%  zweimal bauen: TOC/Kopfzeilen/Tabellennummern (6.x→7.x) und alle
%  Querverweise erneuern sich erst im Rebuild.]

\section{Evaluationsziele und Untersuchungsumfang}
\label{sec:eval-ziele}

Die in Kapitel~\ref{chap:explorative-problemklaerung} konsolidierten
Designanforderungen A1--A10 bilden den Ausgangspunkt der Evaluation.
Fünf Untersuchungsbereiche verbinden die Anforderungen mit konkreten
Prüffragen. Diese Bündelung dient der Prüfung und unterscheidet sich
von den sechs Architektursäulen aus Kapitel~\ref{chap:loesungsarchitektur}.

\emph{Evidenzbindung} (A3) fragt nach der Rückführbarkeit von Claims
auf Quellenbelege und von weiteren fachlichen Inhalten auf ihre
Eingaben oder Entscheidungen. \emph{Kontrollierte Transformation
und Mutation} (A4/A6) untersucht Inhaltserhaltung und Schutzmechanismen
bei getrennten Schritten für Erzeugung, Prüfung, Reparatur und
Autorisierung. \emph{Fortschreibung} (A5/A7) betrifft den Umgang mit
neuen, präzisierenden, widersprechenden, wiederholten und offenen
Informationen. \emph{Projektion} (A8) prüft die Außendarstellung des
Core und die Verarbeitung externer Eingänge; \emph{Durabilität mit
agentischer Bedienung} (A9/A10) betrifft Wiederaufnahme und Steward-Zugang.
Tabelle~\ref{t:eval-pruefkarte} ordnet diesen Bereichen acht Prüffragen
zu. Abschnitt~\ref{sec:eval-ergebnisprofil} führt die Antworten mit
Evidenzform und verbleibender Grenze zusammen.

Die Querschnittsanforderungen A1 (beobachtbare Aktionen) und A2
(explizite Informationsflüsse) erhalten keine eigene Prüffrage.
Ihre Umsetzung wird daran untersucht, welche Verarbeitungsschritte
sich aus Laufartefakten und Ereignisketten rekonstruieren lassen:
die Herkunft der Prüfhinweise in Abschnitt~\ref{sec:eval-luecken},
die Wiederaufnahme in Abschnitt~\ref{sec:eval-resume} und die
Herkunftspfade im Audit in Abschnitt~\ref{sec:eval-provenienz}.

\input{\tabledir/eval-pruefkarte.tex}

Drei Evidenzklassen kehren dabei wieder: kontrollierte
Benchmark-Auswertung gegen einen Referenzbestand [B],
deterministische Audits und Tests [VA] sowie dokumentierte
Demonstrationen und Fallprüfungen [D]. Die Verfahren richten sich nach der zu
prüfenden Eigenschaft: Inhaltserhaltung erfordert einen Vergleich
mit Quellen oder vorherigen Artefaktständen. Für die Durchsetzung
von Freigabe- und Integritätsregeln werden gezielt gültige und
ungültige Fälle geprüft. Fortschreibung und Projektion werden an
Eingaben, Entscheidungen und gespeicherten Zustandsänderungen
nachvollzogen; Ausführungsprotokolle und Prozessgrenzentests
prüfen die Wiederaufnahme. Damit vergleichen die Untersuchungen
die jeweiligen Lösungsziele mit beobachteten Ergebnissen
(Abschnitt~\ref{sec:forschungsdesign}). Ihre Reichweite bleibt
unterschiedlich: Ein Fallbeleg oder bestandener Kontrolltest
begründet keine allgemeine Zuverlässigkeits- oder Überlegenheitsaussage.

Die Untersuchungen finden überwiegend unter künstlichen Prüfbedingungen
statt: an vorgegebenen Quellen, archivierten Fällen und kontrollierten
Tests. Sie bewerten ausgewählte technische und fachliche Eigenschaften;
die Wirksamkeit im laufenden Einsatz durch ein Projektteam wird damit
nicht geprüft. Der Evaluationsrahmen FEDS (Framework for Evaluation
in Design Science) unterscheidet Bewertungszweck und
Untersuchungssituation und unterstützt die Auswahl der Verfahren nach
Zielen und verfügbaren Ressourcen \cite{venable2016feds}.
Dieser Rahmen dient hier der nachträglichen Einordnung der tatsächlich
durchgeführten Untersuchungen.

Zur Reichweite des Konfigurationsvergleichs in Abschnitt~\ref{sec:eval-armf}:
Die Vergleichsbedingung~F verwendet einen Einzelagenten mit
Werkzeugen für Quellenzugriff, deterministische Prüfung und
Ergebniseinreichung. Untersucht werden ausgewählte gemeinsame
Ergebnis- und Rechenschaftseigenschaften der
\emph{Evidenzerzeugung} --- des ersten Kettenglieds. Der Ledger
wird darüber hinaus über seine gespeicherten
Verarbeitungsschritte und seine Übergaben an die Folgestufen
untersucht; die übrigen Systempfade (Kontrollen, Fortschreibung,
Projektion, Wiederaufnahme) haben eigene Nachweise und sind keine
exklusiven Leistungen der Ledger-Stufe. Die Abdeckungsdifferenz
aus Abschnitt~\ref{sec:eval-armf} ist deshalb keine Bewertung der Gesamtarchitektur ---
der Ergebnisvergleich behält innerhalb dieser Messgrenze seine
Aussagekraft.

\textbf{Außerhalb des Untersuchungsumfangs} bleiben ausdrücklich:
Wirtschaftlichkeits-, Zeit- und Produktivitätsmessungen
(berichtet werden nur Token-Indikatoren), Usability der
Bedienschicht, die Güte der Facetten- und Dispositionstaxonomie,
Angriffe außerhalb der Speichernaht (G04), die Qualität der
deterministischen Quellsegmentierung sowie jede Generalisierung
auf unbekannte Quellen, andere Domänen oder andere Modelle über
die im Anhang dokumentierte Modellreihe hinaus.
